```latex
\begin{center}
    \textbf{\Huge{Visão inicial do projeto}}
\end{center}

\section{Introdução}

Os condomínios residenciais e empresariais vêm passando por um processo crescente de modernização, impulsionado pela necessidade de aumentar a segurança, melhorar a comunicação entre moradores e administração e otimizar o gerenciamento das atividades diárias. Nesse contexto, a utilização de ferramentas tecnológicas tornou-se um fator importante para garantir maior eficiência operacional e melhor experiência para todos os envolvidos.

A administração condominial envolve uma grande quantidade de processos e informações que precisam ser acompanhados constantemente, como controle de acesso de visitantes e prestadores de serviço, gerenciamento de encomendas, reservas de áreas comuns, comunicação com moradores, emissão de comunicados, registro de ocorrências e acompanhamento de manutenções.

Entretanto, muitos condomínios ainda realizam parte significativa dessas atividades por meio de controles manuais, utilizando planilhas eletrônicas, aplicativos de mensagens instantâneas, registros físicos e processos descentralizados. Essa forma de gerenciamento pode ocasionar falhas de comunicação, perda de informações, retrabalho, dificuldades de rastreabilidade e baixa eficiência operacional.

Além disso, a falta de integração entre os diferentes setores envolvidos, como portaria, administração, síndico e moradores, dificulta o acesso rápido às informações e torna os processos mais suscetíveis a erros. Em condomínios de médio e grande porte, essas limitações impactam diretamente a segurança, a organização e a qualidade dos serviços prestados aos moradores.

Diante desse cenário, observa-se a necessidade de uma plataforma capaz de centralizar e integrar os principais processos operacionais de um condomínio, permitindo maior controle das atividades, melhoria da comunicação e aumento da eficiência administrativa. Assim, este trabalho propõe o desenvolvimento de um sistema web multiusuário voltado para a gestão condominial, buscando proporcionar segurança, organização, rastreabilidade e transparência para todos os participantes do ambiente condominial.

\section{Posicionamento}

\subsection{Declaração do problema}

\begin{table}[h]
    \begin{tabular}{|p{0.3\textwidth}|p{0.65\textwidth}|}
        \hline
            O problema da & Falta de integração e controle das operações condominiais \\
        \hline
            afeta & Administradores, síndicos, porteiros, funcionários, prestadores de serviço e moradores \\
        \hline
            cujo impacto & A dificuldade de comunicação, o retrabalho operacional, a perda de informações e a baixa rastreabilidade dos processos \\
        \hline
            uma solução bem sucedida & Seria a centralização das informações e processos em uma plataforma integrada de gestão condominial \\
        \hline
    \end{tabular}
\end{table}

\subsection{Descrição da posição do produto}

\begin{table}[h]
    \begin{tabular}{|p{0.2\textwidth}|p{0.75\textwidth}|}
        \hline
            Para & Condomínios residenciais e empresariais \\
        \hline
             O \companyname & é um sistema integrado de gestão condominial \\
        \hline
            Isso & Permitirá o gerenciamento centralizado das atividades operacionais, melhorando a comunicação, a segurança e o controle das informações. \\
        \hline
            Ao contrário & Dos métodos tradicionais que utilizam controles manuais, planilhas e aplicativos de mensagens dispersos. \\
        \hline
            Nosso produto & Proporcionará organização, segurança, rastreabilidade, automação e transparência dos processos. \\
        \hline
    \end{tabular}
\end{table}

\section{Descrições das Partes Interessadas e Usuários}

\subsection{Partes Interessadas}

\begin{table}[h]
    \begin{tabular}{| p{0.2\textwidth} | p{0.3\textwidth} | p{0.42\textwidth} |}
        \hline \textbf{Nome} & \textbf{Descrição} & \textbf{Responsabilidades} \\ \hline

            Administração do condomínio &
            Responsável pela gestão operacional e administrativa do condomínio &
            Gerenciar usuários, processos, ocorrências, reservas e indicadores operacionais
            \\ \hline

            Síndico &
            Responsável pela supervisão geral do condomínio &
            Acompanhar indicadores, aprovar solicitações e fiscalizar operações
            \\ \hline

    \end{tabular}
\end{table}

\newpage

\subsection{Resumo do Usuário}

\begin{table}[h]
    \begin{tabular}{| p{0.30\textwidth} | p{0.62\textwidth} | }
        \hline \textbf{Nome} & \textbf{Descrição} \\ \hline

        Moradores &
        Precisam receber comunicados, acompanhar encomendas, autorizar visitantes, reservar áreas comuns e registrar solicitações
        \\
        \hline

        Porteiros &
        Precisam controlar acessos, registrar visitantes, gerenciar entregas e consultar informações dos moradores
        \\
        \hline

        Síndicos &
        Precisam acompanhar indicadores, ocorrências, comunicados e solicitações dos moradores
        \\
        \hline

        Administradores &
        Precisam gerenciar o condomínio, usuários, reservas, documentos e configurações do sistema
        \\
        \hline

        Prestadores de Serviço &
        Precisam ter seus acessos autorizados e registrados quando realizarem atividades dentro do condomínio
        \\
        \hline

    \end{tabular}
\end{table}

\subsection{Ambiente do Usuário}

Os usuários utilizarão o sistema em diferentes contextos operacionais. Porteiros poderão acessar a plataforma em computadores localizados na portaria, enquanto moradores, síndicos e administradores poderão utilizar dispositivos móveis ou computadores pessoais para acompanhar informações e executar operações relacionadas ao condomínio.

\subsection{Resumo das Principais Partes Interessadas e Necessidades do Usuário}

Atualmente, grande parte das informações é controlada por meios descentralizados, dificultando o acompanhamento das atividades e a comunicação entre os envolvidos. As partes interessadas demonstram necessidade de uma solução que centralize os processos, aumente a segurança e facilite o gerenciamento das operações diárias.

\subsection{Alternativas e Competições}

Existem soluções disponíveis no mercado voltadas para a gestão condominial. Entretanto, muitas apresentam elevado custo de implantação, funcionalidades excessivamente complexas para condomínios de pequeno e médio porte ou limitações de personalização. Dessa forma, existe espaço para uma solução flexível e adaptada às necessidades específicas dos usuários.

\section{Resumo do Produto}

O produto será disponibilizado por meio de uma plataforma web integrada, capaz de centralizar informações, automatizar processos e disponibilizar painéis de acompanhamento para moradores, administração e portaria.

\subsection{Perspectiva do Produto}

O sistema será independente e fornecerá ferramentas para gerenciamento de acessos, visitantes, encomendas, reservas de áreas comuns, ocorrências, comunicados e indicadores operacionais.

\subsection{Suposições e Dependências}

Será necessária uma infraestrutura capaz de hospedar a aplicação e armazenar os dados do condomínio. A comunicação entre os usuários dependerá de conexão com a internet, podendo futuramente ser estudadas alternativas para funcionamento parcial em modo offline.

\section{Outros Requisitos do Produto}

O sistema deverá garantir alta disponibilidade, segurança das informações, proteção de dados pessoais, facilidade de uso e estabilidade operacional, assegurando o funcionamento contínuo das atividades do condomínio.
```
